\documentclass[11pt]{article}

\usepackage{nmrdoc}
\usepackage{siunitx}
\usepackage{bm}

\sisetup{per-mode=symbol}
\DeclareSIUnit{\angstrom}{\text{\AA}}
\DeclareSIUnit{\elementarycharge}{e}
\emergencystretch=12em

\begin{document}
\NmrTitle{The HydrationShell Calculator}

\section{What the calculator computes}

The HydrationShell calculator describes the local arrangement of explicit water
and ions around each protein atom. It does not compute an electric field, an
electric-field gradient, or a shielding tensor. Its output is a four-column
geometric descriptor: a first-shell exposed-water fraction, a mean water dipole
orientation, the distance to the nearest ion within a cutoff, and the charge of
that ion.

These quantities bear on NMR shielding because the electronic response of a
protein atom depends on the local solvent environment. A backbone oxygen facing
ordered water is not in the same electrostatic and hydrogen-bond environment as
one buried against protein atoms, even when the protein coordinates alone look
similar. The calculator therefore supplies solvent-structure variables that can
be calibrated against a DFT shielding reference. Before that calibration, the
numbers are not chemical shifts and not shielding contributions in ppm.

\section{Where shielding enters}

For a nucleus in an applied magnetic field \(\bm B_0\), the local field is
written
\begin{equation}
  \bm B_{\mathrm{nuc}} = (\bm I-\bm\sigma)\bm B_0 .
  \label{eq:shielding}
\end{equation}
Here \(\bm I\) is the \(3\times3\) identity matrix and \(\bm\sigma\) is the
nuclear shielding tensor. The tensor is dimensionless: it is the linear
coefficient that tells how the electrons reduce or redirect the applied field at
the nucleus. A tensor is needed because a molecule has orientation. Rotating the
molecule can change which component of the electron-induced field appears along
which laboratory axis.

The usual solution NMR observable is the isotropic chemical shift, obtained from
the orientational average of \(\bm\sigma\) after reference subtraction and the
chosen sign convention. That scalar average does not exhaust the shielding
information. The anisotropic part records how the shielding changes with
molecular orientation, and it is the part that many geometry calculators in this
pipeline try to preserve as a rank-2 object.

HydrationShell sits beside those rank-2 calculators rather than imitating them.
Its physical claim is weaker and more concrete: nearby solvent geometry carries
information about the dielectric and hydrogen-bond setting of an atom. Water
positions, water dipole directions, and ion contacts can correlate with the DFT
shielding response because they report the boundary condition seen by the local
electrons. The calculator stops at those descriptors. It does not solve the
electronic response of the solvated system.

\section{The tensor split used elsewhere}

When a calculator emits a \(3\times3\) tensor \(\bm X\), the code stores it with
the \texttt{SphericalTensor} split. The scalar part is
\begin{equation}
  T_0 = \frac{1}{3}\operatorname{tr}\bm X ,
  \label{eq:t0}
\end{equation}
the trace divided by three. This is the isotropic component, the part that
survives rotational tumbling.

The antisymmetric part is stored as a three-component pseudovector,
\begin{equation}
  \bm T_1 =
  \frac{1}{2}
  \begin{pmatrix}
    X_{yz}-X_{zy}\\
    X_{zx}-X_{xz}\\
    X_{xy}-X_{yx}
  \end{pmatrix}.
  \label{eq:t1}
\end{equation}
It stores the part of the matrix that changes sign under transpose. This part is
usually not observed directly in standard isotropic solution NMR.

The remaining part is the symmetric traceless matrix
\begin{equation}
  \bm S=\frac{\bm X+\bm X^{\mathsf T}}{2}-T_0\bm I,
  \qquad \operatorname{tr}\bm S=0 .
  \label{eq:s}
\end{equation}
It has five independent entries. The implementation packs those entries as
\begin{equation}
  \bm T_2 =
  \begin{pmatrix}
    \sqrt{2}\,S_{xy}\\
    \sqrt{2}\,S_{yz}\\
    \sqrt{3/2}\,S_{zz}\\
    \sqrt{2}\,S_{xz}\\
    (S_{xx}-S_{yy})/\sqrt{2}
  \end{pmatrix}.
  \label{eq:t2}
\end{equation}
The factors in Eq.~\eqref{eq:t2} preserve length: the squared norm of this
five-vector equals the Frobenius norm \(\sum_{ab}S_{ab}^2\) of the Cartesian
symmetric-traceless matrix. The full packed order is
\[
  [T_0,\ T_{1x},\ T_{1y},\ T_{1z},\
   T_{2,-2},\ T_{2,-1},\ T_{2,0},\ T_{2,+1},\ T_{2,+2}] .
\]

HydrationShell does not call this decomposition. It writes four scalar channels,
not a nine-component tensor and not a five-component \(T_2\) block. The tensor
framework still matters for interpretation of the shielding model: these
scalar solvent descriptors may help calibrate a rank-2 shielding response, but
they are not themselves tensor components.

\section{The method implemented}

For a conformation with \(N\) protein atoms, let \(\bm r_i\) be the position of
atom \(i\) in \(\si{\angstrom}\). The code first forms the protein coordinate
centroid
\begin{equation}
  \bm c = \frac{1}{N}\sum_{j=1}^{N}\bm r_j .
  \label{eq:centroid}
\end{equation}
This is an unweighted average of atom positions, not a mass center. For each
target atom, the direction
\begin{equation}
  \hat{\bm u}_i =
  \frac{\bm c-\bm r_i}{|\bm c-\bm r_i|}
  \label{eq:interior}
\end{equation}
points from the atom toward the protein interior as approximated by that
centroid. The half-shell calculation is a signed divider through the
atom: the plane perpendicular to \(\hat{\bm u}_i\) cuts space into an
interior-facing side and an exterior-facing side. A water oxygen on the centroid
side counts as buried; a water oxygen on the other side counts as exposed.

The first hydration shell for atom \(i\) is the set
\begin{equation}
  \mathcal W_i =
  \{w:\ |\bm o_w-\bm r_i|\le R_{\mathrm w}\},
  \qquad R_{\mathrm w}=\SI{3.5}{\angstrom},
  \label{eq:shell}
\end{equation}
where \(\bm o_w\) is the oxygen position of water molecule \(w\). The cutoff is
\texttt{water\_first\_shell\_cutoff} in \texttt{data/calculator\_params.toml}.
Hydrogen positions do not affect shell membership.

For each first-shell water, define the unit vector from the atom to the water
oxygen,
\begin{equation}
  \hat{\bm r}_{iw} =
  \frac{\bm o_w-\bm r_i}{|\bm o_w-\bm r_i|}.
  \label{eq:rhat}
\end{equation}
The code classifies the water by the sign of
\(\hat{\bm r}_{iw}\cdot\hat{\bm u}_i\):
\begin{align}
  n_{\mathrm{buried},i} &=
  \sum_{w\in\mathcal W_i}
  \mathbf 1(\hat{\bm r}_{iw}\cdot\hat{\bm u}_i > 0),\\
  n_{\mathrm{exposed},i} &=
  \sum_{w\in\mathcal W_i}
  \mathbf 1(\hat{\bm r}_{iw}\cdot\hat{\bm u}_i \le 0).
  \label{eq:counts}
\end{align}
The equality case goes to the exposed side because the source uses
\texttt{> 0} for buried and \texttt{else} for exposed. The written
half-shell feature is
\begin{equation}
  h_i =
  \begin{cases}
    \dfrac{n_{\mathrm{exposed},i}}
          {n_{\mathrm{exposed},i}+n_{\mathrm{buried},i}},
       & n_{\mathrm{exposed},i}+n_{\mathrm{buried},i}>0,\\[1.1em]
    0, & n_{\mathrm{exposed},i}+n_{\mathrm{buried},i}=0 .
  \end{cases}
  \label{eq:half-shell}
\end{equation}
This value is a fraction in \([0,1]\). It is not a signed asymmetry about zero,
and HydrationShell does not write the raw first-shell count. Thus \(h_i=0\) can
mean either that the first-shell waters are on the buried side or that no
first-shell water was found.

The orientation channel uses the water dipole vector stored by
\texttt{SolventEnvironment}. With oxygen as the origin of the water molecule,
the source computes
\begin{equation}
  \bm\mu_w =
    q_{\mathrm H}(\bm h_{1w}-\bm o_w)
  + q_{\mathrm H}(\bm h_{2w}-\bm o_w),
  \label{eq:dipole}
\end{equation}
where \(\bm h_{1w}\) and \(\bm h_{2w}\) are hydrogen positions and
\(q_{\mathrm H}\) is the hydrogen charge read from the GROMACS topology, in
\(\si{\elementarycharge}\). The oxygen charge does not appear because its
displacement from the chosen origin is zero. The vector \(\bm\mu_w\) has units
\(\si{\elementarycharge\angstrom}\). Waters with
\(|\bm\mu_w|\le 10^{-10}\) in those units are skipped; this is the
\texttt{near\_zero\_vector\_norm\_threshold}.

For the remaining first-shell waters, the code averages
\begin{equation}
  d_i =
  \frac{1}{m_i}
  \sum_{\substack{w\in\mathcal W_i\\ |\bm\mu_w|>10^{-10}}}
  \frac{\bm\mu_w\cdot\hat{\bm r}_{iw}}{|\bm\mu_w|},
  \label{eq:dipole-cos}
\end{equation}
where \(m_i\) is the number of waters included in the sum. If \(m_i=0\), the
written value is \(0\). The term inside the sum is a cosine. With the usual
positive hydrogen charge, a positive value means the water dipole points in the
same direction as the atom-to-oxygen vector; geometrically, the hydrogens lie on
average farther from the atom than the oxygen. A negative value means the
hydrogen side points more toward the atom.

Ion proximity is computed independently of the water shell. When at least one
ion is present, with ion positions \(\bm z_k\) and charges \(q_k\), the code
finds
\begin{equation}
  k_i^\star = \arg\min_k |\bm z_k-\bm r_i|,
  \qquad
  \rho_i = |\bm z_{k_i^\star}-\bm r_i|.
  \label{eq:nearest-ion}
\end{equation}
It then writes
\begin{equation}
  I_i =
  \begin{cases}
    \rho_i, & \rho_i < R_{\mathrm{ion}},\\
    +\infty, & \rho_i \ge R_{\mathrm{ion}},
  \end{cases}
  \qquad
  Q_i =
  \begin{cases}
    q_{k_i^\star}, & \rho_i < R_{\mathrm{ion}},\\
    0, & \rho_i \ge R_{\mathrm{ion}},
  \end{cases}
  \label{eq:ion-output}
\end{equation}
with \(R_{\mathrm{ion}}=\SI{20.0}{\angstrom}\), the
\texttt{hydration\_ion\_cutoff}. The comparison is strict: an ion exactly at the
cutoff is reported as absent. The ion distance has units
\(\si{\angstrom}\); the ion charge has units
\(\si{\elementarycharge}\). The charge channel's zero should be interpreted with
the distance channel, since \(Q_i=0\) accompanies ``no ion in cutoff''.
If the solvent environment contains no ions, the same absent branch writes
\(+\infty\) for \(I_i\) and \(0\) for \(Q_i\).

The implementation loops over each water and each ion for each atom. Although
\texttt{SpatialIndexResult} is a declared dependency, this routine does not use
a spatial index in its computation. The trajectory reader supplies coordinates
in \(\si{\angstrom}\), converted from GROMACS nanometres. HydrationShell receives
no box matrix, so the distances above are direct coordinate differences after
the protein has been made whole; no minimum-image solvent distance appears in
this calculator.

The main approximations follow from those definitions. The shell boundary is a
fixed \(\SI{3.5}{\angstrom}\) oxygen cutoff, not a frame-specific radial
distribution minimum. Exposure is measured by a centroid direction, not by a
solvent-accessible surface normal. A separate calculator uses the surface-normal
reference frame; HydrationShell keeps the centroid construction. The model also
assumes that no target atom lies exactly at the coordinate centroid and that no
selected water oxygen is exactly coincident with the target atom, because
Eqs.~\eqref{eq:interior} and \eqref{eq:rhat} would otherwise be undefined.

\section{Inputs and output}

HydrationShell requires a \texttt{SolventEnvironment}. In the trajectory path,
that object is built from a full-system GROMACS TPR and TRR: waters are
identified as three-site O--H--H molecules, ions as single-atom molecules, and
coordinates are split into protein and solvent arrays. The water hydrogen charge
used in Eq.~\eqref{eq:dipole} and the ion charges come from the TPR. Four-site
water models are not handled by the reader used here.

The calculator does not consume MOPAC charges or bond orders, AIMNet2 charges,
ff14SB protein partial charges, APBS potentials, or DSSP assignments. Those
external tools feed other calculators in the same pipeline. HydrationShell's
external input is explicit solvent geometry and topology charge metadata from
the molecular-dynamics frame.

The feature file is \texttt{hydration\_shell.npy} with shape \((N,4)\) and
columns
\[
  [h_i,\ d_i,\ I_i,\ Q_i].
\]
These are scalar descriptors of solvent structure around atom \(i\). Calibration
against a DFT reference is what can turn their fitted combination with other
features into a shielding contribution. Before calibration, the file reports
geometry: exposed-side water fraction, mean first-shell water dipole cosine,
nearest-ion distance, and nearest-ion charge.

\end{document}
